\documentclass[11pt,letterpaper]{article}
\usepackage[margin=.7in,top=.78in,bottom=.7in]{geometry}
\usepackage[T1]{fontenc}
\usepackage{lmodern,amsmath,amssymb,array,booktabs,fancyhdr,hyperref}
\hypersetup{hidelinks}
\renewcommand{\familydefault}{\sfdefault}
\pagestyle{fancy}\fancyhf{}
\fancyhead[L]{\small Week 73 / The gentlest stretch / Adult guide}
\fancyfoot[L]{\small Bellingham Math Circle / Week 73 / GGT73-FAC-v1}
\fancyfoot[R]{\small\thepage}
\renewcommand{\headrulewidth}{0pt}\renewcommand{\footrulewidth}{0pt}
\setlength{\parindent}{0pt}\setlength{\parskip}{6pt}
\setlength{\headheight}{14pt}
\newcommand{\sect}[1]{\par\vspace{5pt}{\large\bfseries #1}\par}
\newcommand{\prob}[1]{\par\vspace{3pt}\textbf{Problem #1}\par}
\begin{document}

\sect{The mathematics first}
For Euclidean straight-line distances, the least possible worst stretch from a
side-marked $4$-by-$1$ rectangle onto a $2$-by-$2$ square is \textbf{2}.
Here a legal whole-sheet map is a homeomorphism: continuous, one-to-one and onto,
with a continuous inverse, taking each named side to its matching side.
The worst stretch means the supremum of new distance divided by old distance
across all distinct pairs. Maps with unbounded stretch cannot improve the answer.

Two facts do the work. A vertical unit segment joins the old bottom and top;
its images are at least $2$ apart. Conversely, halving every horizontal distance
and doubling every vertical distance gives a legal map whose every pair stretches
by at most $2$, including diagonal pairs. For a target of width $W>0$ and height
$H>0$, the same argument gives the exact answer $\max(W/4,H)$.

The opening five-dot game deliberately hides information: if all ten pairs are
checked, \emph{every} interior choice of the new O scores $2$. In the four-triangle
affine model, two extra side-midpoint tests single out the centered choice exactly.
Children may discover that a good test score is not a proof about every point.
Keep three stages distinct: observed ruler results; an exact argument about this
fan family; a bound for every pair and every legal whole-sheet map.

\textbf{Who this fits.} One Grades 4--5 prototype, for children able to compare
lengths, use halves and ratios, follow matching points and coordinate partner roles.
An adult can read aloud or record. The midpoint-disk and diagonal proofs are
readiness-dependent extensions; squared-distance algebra is an adult resource.
There is no separate K--1 or Grades 2--3 route and no tested age-fit claim.

\sect{Materials and preparation}
For the older table of three children, use a pair plus a rotating tester/referee
and one anchored adult. Per pair: one four-page student packet, two or three spare
copies of pages 1--2, five small labeled paper dots, pencil, ruler, scrap paper
strips and blank grid paper. Keep O movable; fixed corners can remain printed.
Print at actual size, without ``fit to page.'' On pages 1--3, one drawn unit is
$0.94$ inches; the quarter-unit grid helps locate points. Page 4 uses its own
shared scale within each old/new comparison. Never compare a length from one
scale with a length from another.

Before use, rehearse placing a dot, marking strip lengths and locating a side's
midpoint. No physical rehearsal has yet occurred. The intended measurement task
must tolerate small errors: even an O only a quarter-unit sideways from center
produces about $0.735$ mm excess in the midpoint test at actual size. Still smaller
offsets can be undetectable by ruler.

\sect{Launch and a feasible first visit}
Give children a minute to handle strips and dots. Demonstrate the printed
$2$ cm to $3$ cm example: the ratio is $3/2=1\frac12$. Show one legal O placement
strictly inside the square and one pair test, without selecting the revealing
midpoints. Let partners choose placements and tests, then exchange jobs.

Allow roughly 10--15 minutes for Problem 1. Move on when the apparent tie becomes
interesting; repeated placements are not the goal. Spend the next 15--20 minutes
on Problem 2, showing its midpoint convention only when needed. A good stopping
point is ``our first tests missed something'' plus a candidate centered map.
Use Problems 3--4 as a later discussion or continuation; retain Problems 5--6
for children ready to invent and justify whole-sheet rules.
\newpage
\sect{Answers and optional hints}
Use normalized coordinates with the old lower-left corner at $(0,0)$ and the new
one at $(0,0)$; the page's vertical offset is just drawing layout.

\prob{1}
\textbf{Every interior O ties at 2 if all original pairs are checked.}
The ten pairs consist of six corner pairs and four O-to-corner pairs.
Old vertical corner pairs have lengths $1\to2$, so their stretch is $2$.
Horizontal corner pairs have $4\to2$, so their stretch is $1/2$.
Corner diagonals have $\sqrt{17}\to\sqrt8$, a ratio below $2$.
Old O is $(2,1/2)$, at distance $\sqrt{17}/2$ from each corner.
New interior O is less than $\sqrt8$ from every corner; these four ratios are
less than $2$ as well. Testing only some pairs can give a smaller \emph{reported}
score, which is why the partner hunts for the largest stretch found.

\textit{Hint if needed:} ``Which pair is unchanged when O moves?''
Do not announce that every placement ties before children have tried contrasting
ones. This intentional failure of sparse testing prepares Problem 2.

\prob{2}
\textbf{Exactly the center passes the two decisive midpoint tests mathematically.}
Let M and N be the midpoints of the bottom and top sides. In the old sheet,
$OM=ON=1/2$; in the new sheet $M'=(1,0)$ and $N'=(1,2)$.
Thus an O image $(u,v)$ can pass both stretch-at-most-$2$ tests only if
\[
 (u-1)^2+v^2\le1,\qquad (u-1)^2+(v-2)^2\le1.
\]
Adding and simplifying gives $(u-1)^2+(v-1)^2\le0$, so $O'=(1,1)$.
A geometric version draws the two closed radius-$1$ disks centered at M' and N';
they meet only at their common midpoint. Every off-center O lies outside at
least one disk. The centered fan passes every pair by the map in Problem 3.

The model itself is legal: O must be strictly interior. Each old triangle maps
affinely onto its matching new triangle; matching edge maps agree. The triangles
partition each sheet, so gluing them gives a side-preserving homeomorphism.
A halfway point on an edge maps to the corresponding halfway point, as shown
in the student example. \textit{Hint:} offer side midpoints after children have
tried their own added dots.

\textbf{Measurement limit.} The printed task asks which positions survive
\emph{their tests}. Accept an honest set of apparently surviving near-center
positions. Rulers cannot certify the unique exact point or arbitrarily small
excess stretch. The tangent-disk reasoning, rather than finer measuring, settles
that exact question.

\prob{3}
\textbf{One valid whole-sheet rule is $F(x,y)=(x/2,2y)$.}
Children may describe it as halving distance across the sheet and doubling height.
Old dotted vertical lines spaced $1/2$ unit apart become lines $1/4$ unit apart;
the old middle horizontal line becomes the square's middle horizontal line.
The map fills the square without gaps or overlaps and preserves every named side.
Its inverse doubles across and halves height.

\prob{4}
\textbf{Yes for that rule, and no rule can make all ratios less than 2.}
For any two points whose coordinate differences are $(p,q)$,
\[
 |F(P)-F(Q)|^2=p^2/4+4q^2\le4(p^2+q^2)=4|P-Q|^2.
\]
This covers diagonal pairs and points between grid lines. Geometrically, first
double the sheet uniformly, then compress horizontally; the latter operation
cannot increase a straight-line distance. For the lower bound, use an old
vertical unit pair on the two horizontal boundaries. Every legal map sends its
endpoints to opposite square sides, at distance at least $2$.
\newpage
\sect{Changing the target}
\prob{5}
\textbf{The $6$-by-$2$ target has optimum 2.} Use
$F(x,y)=(3x/2,2y)$. The vertical unit pair still forces ratio at least $2$,
and both coordinate factors are at most $2$.

\textbf{The $2$-by-$3$ target has optimum 3.} Use
$F(x,y)=(x/2,3y)$. An old vertical unit pair joins target sides at least $3$
apart, and neither coordinate factor exceeds $3$.

For either rule, a squared-distance argument like Problem 4 handles every pair;
success on the displayed grid alone is not enough. The unavoidable pair is an
old bottom-to-top pair with the same horizontal coordinate. Its image need not
be vertical under an arbitrary map, but its vertical separation alone gives
the lower bound. \textit{Hint:} ``How far apart must the new top and bottom be?''

\prob{6}
\textbf{Examples include $6$ by $1$ and $3$ by $1\frac12$.}
Their rules are $(3x/2,y)$ and $(3x/4,3y/2)$ respectively. Both have optimum
$1\frac12$, with a horizontal boundary pair forcing the first answer and a
vertical boundary pair forcing the second.

The complete answer is all positive widths and heights satisfying
\[
 W\le6,\qquad H\le1\tfrac12,\qquad
 \text{at least one equality.}
\]
Indeed, any legal map has worst stretch at least $W/4$, by an old horizontal
length-$4$ pair on opposite vertical sides, and at least $H$, by an old vertical
unit pair. The diagonal linear map $(x,y)\mapsto(Wx/4,Hy)$ attains
$\max(W/4,H)$: with $L=\max(W/4,H)$,
\[
 (Wp/4)^2+(Hq)^2\le L^2(p^2+q^2).
\]
This also excludes invented rectangles for which both factors are smaller than
$1\frac12$, or one is larger. \textit{Optional extension:} invite an example
where the width is the limiting constraint and another where the height is.

\sect{What is proved and what is still unknown}
The boundary and squared-distance arguments prove the all-maps optimum under
the stated side correspondence. The two-disk argument applies specifically to
the affine fan family. It does not classify all optimal homeomorphisms.
Finite pin tests only certify the pairs actually tested. An arbitrary sketch
without a whole-sheet rule cannot establish an upper bound.

This is a one-sided Euclidean Lipschitz optimization. Do not identify its number
with a Teichm\"uller distance or a quasiconformal normalization, and do not drop
the side-preserving hypothesis. Real rubber need not stretch uniformly; it is
not a proof model for this packet.

\sect{Source lineage and status}
Metric-comparison inspiration: Moon Duchin, \emph{Billiards and Geometric
Topology}, Reed colloquium, November 3, 2005,
\href{https://people.reed.edu/~davidp/talks/talks05-06/talks/duchin.html}{talk abstract}.
The abstract discusses comparing metrics and flat structures. The rectangle
optimization, five-pin experiment and proofs here are independently developed;
no claim is made that they appear in that talk.

\textbf{Review status:} matched to the final four-page GGT73-S-v1 student packet.
The guide's independent standard-library checker verifies the printed numerical
instances; universal claims use the proofs above. The PDF has been built,
rendered and visually checked. This is an unpiloted review prototype;
physical setup, ruler usability and classroom pacing remain unrehearsed.

\end{document}
